%%%%%%%%%%%%%%%%%%%%%%%%%%%%%%%%%%%%%%%%%
% Beamer Presentation
% LaTeX Template
% Version 1.0 (10/11/12)
%
% This template has been downloaded from:
% http://www.LaTeXTemplates.com
%
% License:
% CC BY-NC-SA 3.0 (http://creativecommons.org/licenses/by-nc-sa/3.0/)
%
%%%%%%%%%%%%%%%%%%%%%%%%%%%%%%%%%%%%%%%%%

%----------------------------------------------------------------------------------------
%	PACKAGES AND THEMES
%----------------------------------------------------------------------------------------

\mode<presentation> {

% The Beamer class comes with a number of default slide themes
% which change the colors and layouts of slides. Below this is a list
% of all the themes, uncomment each in turn to see what they look like.

%\usetheme{default}
%\usetheme{AnnArbor}
%\usetheme{Antibes}
%\usetheme{Bergen}
%\usetheme{Berkeley}
%\usetheme{Berlin}
%\usetheme{Boadilla}
%\usetheme{CambridgeUS}
%\usetheme{Copenhagen}
%\usetheme{Darmstadt}
%\usetheme{Dresden}
%\usetheme{Frankfurt}
%\usetheme{Goettingen}
%\usetheme{Hannover}
%\usetheme{Ilmenau}
%\usetheme{JuanLesPins}
%\usetheme{Luebeck}
%\usetheme{Madrid}
%\usetheme{Malmoe}
%\usetheme{Marburg}
%\usetheme{Montpellier}
%\usetheme{PaloAlto}
%\usetheme{Pittsburgh}
%\usetheme{Rochester}
%\usetheme{Singapore}
%\usetheme{Szeged}
%\usetheme{Warsaw}

% As well as themes, the Beamer class has a number of color themes
% for any slide theme. Uncomment each of these in turn to see how it
% changes the colors of your current slide theme.

%\usecolortheme{albatross}
%\usecolortheme{beaver}
%\usecolortheme{beetle}
%\usecolortheme{crane}
%\usecolortheme{dolphin}
%\usecolortheme{dove}
%\usecolortheme{fly}
%\usecolortheme{lily}
%\usecolortheme{orchid}
%\usecolortheme{rose}
%\usecolortheme{seagull}
%\usecolortheme{seahorse}
%\usecolortheme{whale}
%\usecolortheme{wolverine}

% To remove the footer line in all slides uncomment this line
%\setbeamertemplate{footline}
% To replace the footer line in all slides with a simple slide count uncomment
% this line
\setbeamertemplate{footline}[page number]
% To remove the navigation symbols from the bottom of all slides uncomment
% this line
\setbeamertemplate{navigation symbols}{}
}

% Beamer font style
\setbeamerfont{title}{size=\huge,series=\bf}
\setbeamerfont{section title}{size=\huge,series=\bf}
\setbeamerfont{frametitle}{size=\LARGE}
\setbeamerfont{itemize/enumerate body}{size=\large}
\setbeamerfont{normal text}{size=\large}

\makeatother{}

% Input encoding
\usepackage[utf8]{inputenc}
% Fonts and symbols
\usepackage{amsfonts, amsmath, amssymb, amsthm}
\usepackage{eurosym}
% Image
\usepackage{graphicx}
% Colors
\usepackage{color}
% Allows the use of \toprule, \midrule and \bottomrule in tables
\usepackage{booktabs}
% Extended tables
\usepackage{tabularx}
% ITE package (for conditional compilation)
\usepackage{ifthen}

% Color for emphasis
\definecolor{emphasis}{rgb}{0.1,0.1,0.8}

% Listings (code)
\usepackage{listings}
\usepackage{courier}
\lstset{basicstyle=\ttfamily\footnotesize
	,tabsize=2
	%,frame=single
	%,numbers=left
	,numberstyle=\tiny
	,numbersep=5pt
	,numberfirstline=true
	,firstnumber=1
	,stepnumber=5
	,showspaces=false
	,showstringspaces=false
	,showtabs=false
	,breaklines=false
	,breakatwhitespace=true
}
\lstdefinestyle{highlight}{
	keywordstyle=\color{red},
	commentstyle=\color{green},
}
\lstdefinestyle{beamer}{
 	%preset=\large,
    breaklines=true,
	basicstyle=\ttfamily\footnotesize\color{black!50},
	keywordstyle=\color{red!50},
	commentstyle=\color{green!50},
	moredelim=**[is][{\color{emphasis}\lstset{style=highlight}}]{@}{@},
}

% Macro to include images
\newcommand{\includeimage}[2][scale=1.0]{
	\begin{center}
	\includegraphics[keepaspectratio=true,#1]{img/#2}
	\end{center}
}

% Macro: separation between bad and good examples in attacks
\newcommand{\badgoodsep}{
\vspace{6mm}
\pause
\hrule
\vspace{6mm}
}

% Macro for good and bad code
\newenvironment{embeddedcode}[1]{
\begin{columns}[T]
\begin{column}{3mm}
\includeimage[width=5mm]{code-#1}
\end{column}
\begin{column}{0.9\linewidth}
}{
\end{column}
\end{columns}
}

% Macro frametransition
\newenvironment{frametransition}[1]{
\begin{frame}
\begin{center}
\Huge\bfseries #1
\end{center}
}{
\note{
	Transition...
}
\end{frame}
}

% Begin each section with a dedicated frame
\AtBeginSection[]{
	\begin{frame}
	%\sectionpage
	\usebeamercolor[fg]{title}
	\usebeamerfont{section title}
	\centering
	\insertsection
	\note{\Large New section: \emph{\insertsection}}
	\end{frame}
}


%-----------------------------------------------------------------------------
%	TITLE PAGE
%-----------------------------------------------------------------------------

\title[Web Security]{Security in Web Applications}
\author[]{Fran\c{c}ois Roland
\and Krzysztof Magusiak
\and Vincent Gribomont
\and Micha\"{e}l Lacroix
}
\institute[Hermes ECS]{Hermes Engineering Consultancy Services}
\date{\today}

\begin{document}

%-----------------------------------------------------------------------------

\begin{frame}
\author{}
\institute{}
\date{}
\titlepage
% Insert the logo
\begin{picture}(0,0)
\put(160,0){\includegraphics[keepaspectratio,width=5cm]{img/hermes}}
\end{picture}
\end{frame}

%\begin{frame}[allowframebreaks]
%\frametitle{Overview}
%\setcounter{tocdepth}{2}
%\tableofcontents
%\end{frame}
% For the rest of the document set depth to 3
\setcounter{tocdepth}{3}

%-----------------------------------------------------------------------------
%	PRESENTATION SLIDES
%-----------------------------------------------------------------------------

\begin{frame}
\frametitle{Hermes}
\includegraphics[keepaspectratio,width=0.4\linewidth]{img/hermes}
\vspace{1cm}\par
description
\end{frame}

\begin{frame}
\frametitle{François Roland}
\begin{columns}
\begin{column}{0.4\linewidth}
%\includeimage[width=\linewidth]{photo-francois-roland}
\end{column}
\begin{column}{0.6\linewidth}
description
\end{column}
\end{columns}
\end{frame}

\begin{frame}
\frametitle{Krzysztof Magusiak}
\end{frame}

\begin{frame}
\frametitle{Michaël Lacroix}
\end{frame}

\begin{frame}
\frametitle{Vincent Gribomont}
\end{frame}



\section{Psycho-sociological Aspect}

\subsection{Security concerns}

\begin{frame}
\frametitle{Internet Security History}
Army, Universities, ...
\pause\par
\emph{Security?}
\pause
\includeimage[width=\linewidth,height=0.5\textheight]{mainframe}
\note{
The web:
\begin{itemize}
\item Designed in a \emph{trusted environment}
\item Became global
\end{itemize}
Security was not a concern in a closed environment.
}
\end{frame}

\begin{frame}
\frametitle{Nowadays...}
\begin{columns}
\begin{column}{0.3\linewidth}
\includeimage[width=\linewidth]{logo-facebook}
\end{column}
\begin{column}{0.4\linewidth}
\includeimage[width=\linewidth]{visa}
\end{column}
\begin{column}{0.3\linewidth}
\includeimage[width=\linewidth]{logo-tax-on-web}
\end{column}
\end{columns}
\begin{center}
{\Huge It's \textbf{everywhere}!}
\end{center}
\note{
\textbf{Web services} \\
Security concerns various aspects of our internet life.
\begin{itemize}
\item Personal (e-mail, Facebook, ...)
\item Financial (online banking, Tax-On-Web)
\item Professional (e-mail, calendars, ...)
\item etc.
\end{itemize}
}
\end{frame}

\subsection{Sentiment of security}

\begin{frame}
\begin{columns}
\begin{column}{0.45\linewidth}
\includeimage[width=0.9\linewidth]{https}
\end{column}
\begin{column}{0.1\linewidth}
{\centering\Huge $\nRightarrow$}
\end{column}
\begin{column}{0.45\linewidth}
\includeimage[width=0.9\linewidth]{secure-website}
\end{column}
\end{columns}
\note{
\textbf{False sentiment of security} \\
Gain some level of trust...
\\ but, do not lure the users in a false sentiment of security.
\begin{exampleblock}{HTTPS}
Users may think that an HTTPS web site is secured.
\\ HTTPS prevents only some of the attacks.
\end{exampleblock}
}
\end{frame}

\begin{frame}
\begin{columns}
\begin{column}{0.3\linewidth}
\includeimage[width=\linewidth]{logo-psn}
\end{column}
\begin{column}{0.6\linewidth}
\begin{exampleblock}{PlayStation Network outage}
77 million accounts affected.
\end{exampleblock}
\end{column}
\end{columns}
\note{
A web site must earn the confidence of its users.
\begin{block}{Why?}
Users will send their data to the server.
\end{block}
\begin{block}{However...}
The web site could be compromised.
\end{block}
\begin{exampleblock}{PlayStation Network outage}
Information from approximately 77 million accounts was stolen.
\end{exampleblock}
}
\end{frame}

\subsection{Web payment}

\begin{frame}
\includeimage[width=0.8\linewidth]{secure-payment-logos}
\begin{center}
\begin{minipage}{0.5\linewidth} % to get the list centered
\begin{itemize}
\item Security measures
\item Guarantees
\item Protection
\end{itemize}
\end{minipage}
\end{center}
\note{
\begin{center}
You don't trust everyone with your money.
\end{center}
\begin{itemize}
\item Proper security measures
\item Guarantees
\item Protection against fraud
\end{itemize}
\begin{example}
Visa, MasterCard, PayPal, etc.
\end{example}
}
\end{frame}

\subsection{Users as a security flaw}

\begin{frame}
\frametitle{Security}
\includeimage[height=0.75\textheight,width=\linewidth]{xkcd538_security-1}
\note{
How can we hack that security?
}
\end{frame}

\begin{frame}
\frametitle{Security}
\includeimage[height=0.75\textheight,width=\linewidth]{xkcd538_security-2}
\note{
\begin{itemize}
\item Users can divulge their credentials
\item They may use unsecured web sites, compromised software, ...
\item Users don't always understand what they do
\item "Security is the IT job"
\item etc...
\end{itemize}
\begin{block}{Social engineering}
\begin{quote}
We can pay him on some social network to collect information from him
and eventually steal his identity.
\end{quote}
\end{block}
}
\end{frame}

\begin{frame}
\frametitle{Most used passwords}
\begin{scriptsize}
\begin{enumerate}
\item password
\item 123546
\item 12345678
\item abc123
\item qwerty
\item monkey
\item letmein
\item dragon
\item 111111
\item baseball
\item ...
\end{enumerate}
\end{scriptsize}
% TODO fill the blank
\note{
	Remembering multiple passwords is... \emph{difficult}.
	\begin{itemize}
	\item How many different passwords do you have? More than 5?
	\item Is it better to have a larger alphabet or a longer password?
	\end{itemize}
	Most used passwords source:
	\url{http://www.cbsnews.com/8301-205_162-57539366/the-25-most-common-passwords-of-2012/}
}
\end{frame}

\begin{frame}
\includeimage[height=0.75\textheight,width=\linewidth]{password-post-its}
% TODO add URLs to online services on the page
\note{
\textbf{Bad password storage}
\par
... vs tracking all your different passwords.
\par
Online services for storing passwords:
\begin{itemize}
\item \href{https://agilebits.com/onepassword}{1password}
\item \href{https://lastpass.com/index.php?fromwebsite=1}{lastpass}
\item \href{https://www.pwdhash.com/}{pwdhash}
\end{itemize}
}
\end{frame}

\begin{frame}
\begin{columns}
\begin{column}{0.4\linewidth}
\Large
\centering
\emph{password}
\end{column}
\begin{column}{0.2\linewidth}
\includeimage[height=8mm]{vs}
\end{column}
\begin{column}{0.4\linewidth}
\includeimage[width=\linewidth,height=3cm]{pen}
\end{column}
\pause
\end{columns}
\begin{exampleblock}{Infosecurity Europe 2003 survey}
\begin{quote}
In a 2003 information security survey, {\Large 90\%} of office workers gave
researchers what they claimed was their {\Large password} in answer to a
survey question in exchange for a {\Large cheap pen}.
\end{quote}
\end{exampleblock}
\note{
	\begin{exampleblock}{What do you trade your password for?}
	\begin{quote}
	90\% of people gave away their passwords in exchange for a pen.
	--~Infosecurity~Europe~2003
	\end{quote}
	\end{exampleblock}
}
\end{frame}

\subsection{Trade-off: security vs usability}

\begin{frame}
\begin{center}
\begin{tabular}{rcl}
\pause
secure password & \includegraphics[width=5mm]{img/vs} & simple password
\\[2em]
\pause
authorization checks & \includegraphics[width=5mm]{img/vs} & no authentication
\\[2em]
\pause
encoded data & \includegraphics[width=5mm]{img/vs} & plain text
\end{tabular}
\end{center}
\note{
\begin{itemize}
\item The system should be easy to use
\item Security measures are in the way of what the user wants to achieve
\end{itemize}
Security measures:
\begin{itemize}
\item Too many: make a web site unusable
\item Too few: the web site can be compromised
\end{itemize}
}
\end{frame}



\section{Authentication and Identity}

%------------------------------------------------

\begin{frame}
\begin{columns}
\begin{column}{0.5\linewidth}
\begin{block}{Identity}
Who you are.
\end{block}
\end{column}
\begin{column}{0.5\linewidth}
\includeimage[width=\linewidth,height=0.25\textheight]{identity}
\end{column}
\end{columns}
\pause
\begin{columns}
\begin{column}{0.5\linewidth}
\begin{block}{Authentication}
Are you who you claim to be?
\end{block}
\end{column}
\begin{column}{0.5\linewidth}
\includeimage[width=\linewidth,height=0.25\textheight]{authentication}
\end{column}
\end{columns}
\pause
\begin{columns}
\begin{column}{0.5\linewidth}
\begin{block}{Authorization}
What are you allowed to do?
\end{block}
\end{column}
\begin{column}{0.5\linewidth}
\includeimage[width=\linewidth,height=0.25\textheight]{authorization}
\end{column}
\end{columns}

\note{
\begin{block}{Identity}
The identity of somebody is who they claim.
\newline The definition also applies to non-humans (systems).
\end{block}
\begin{block}{Authentication}
Authentication is the act of asserting the identity of someone.
\end{block}
\begin{block}{Authorization}
Authorization is determining whether the user has the permission to do some
action.
\end{block}
}
\end{frame}

%------------------------------------------------
\begin{frame}
\frametitle{Why authentication?}
\begin{center}
\Large{\textbf{Identification and confidentiality}}
\end{center}
\note{
Authentication is needed when you want to transmit
information and be sure that your correspondent isn't
impersonated by somebody else.
\begin{block}{Confidentiality}
Authentication by itself doesn't imply confidentiality.
For that, we have encryption.
\end{block}
}
\end{frame}

%%%%%%%%%%%%%%%%%%%%%%%%%%%%%%%%%%%%%%%%%%%%%%%%%%%%%%%%%%%%%%%%%%%%%%%%%%%%%%
\subsection{Technical introduction}

\subsubsection{HTTP}

\defverbatim[colored]\Lst{
\begin{lstlisting}[style=beamer]
@GET@ @/wiki/Hypertext_Transfer_Protocol@ HTTP/1.1
Host: en.wikipedia.org
User-Agent: Mozilla/5.0 (X11; Linux x86_64; rv:19.0) Gecko/20100101 Firefox/19.0
Accept: text/html,application/xhtml+xml,application/xml;q=0.9,*/*;q=0.8
@Cookie@: @mediaWiki.user.id=m5MFt3bI4Lr1GcMvj4rPk4mwtwZqTWhx@; 
Connection: keep-alive
...
\end{lstlisting}
}

%TODO : transition + display code properly
\begin{frame}
\frametitle{How does HTTP work?}
\Lst
\note{
A \emph{request} example.
\begin{itemize}
\item Method: \texttt{GET}, \texttt{POST}, etc.
\item URL: resource path
\item Options: \texttt{Host}, \texttt{Cookie}, etc.
\end{itemize}
The \emph{response} has the same structure.
\begin{block}{Anonymity}
No identifier for a web user.
\par\textbf{Cookies}\\
Data stored on the client's computer.
\par\textbf{Session number}\\
The server identifies its users using a session ID (unique number) for each
client. But... the client can send anything.
\end{block}
}
\end{frame}

\subsubsection*{Cryptography introduction}


\begin{frame}
\frametitle{Symmetric cryptography}
\includeimage[width=0.8\linewidth]{cryptography-symm}
\note{
\begin{block}{Encryption}
Transformation of a message in a way that it is readable only by the
authorized parties.
\end{block}
\begin{block}{Symmetric}
A single key used to encode and decode the messages.
\[ msg \xrightarrow{key} cipher \xrightarrow{key} msg \]
\end{block}
\begin{exampleblock}{Symmetric cryptography algorithms}
AES, DES,...
\end{exampleblock}
}
\end{frame}

\begin{frame}
\frametitle{Asymmetric cryptography}
\includeimage[width=0.8\linewidth]{cryptography-asym}
\note{
\begin{block}{Asymmetric}
A private-public key pair, one is used to encrypt the message, the other to
read it.
\[ msg \xrightarrow{key_{pub}} cipher \xrightarrow{key_{priv}} msg \]
\end{block}
\begin{exampleblock}{Asymmetric cryptography algorithms}
RSA, DSA, ...
\end{exampleblock}
}
\end{frame}

\begin{frame}
\frametitle{Public-private keys}
\includeimage[width=\linewidth,height=\textheight]{ppk}
\note{
By definition, the public key is shared.
\begin{itemize}
\item Encrypt with the recipient's public key
\item Sign with your private key
\end{itemize}
}
\end{frame}

\begin{frame}
\begin{center}
\begin{large}
\fontsize{16pt}{25}\selectfont
\begin{tabular}{c|c}
\bfseries Symmetric & \bfseries Asymmetric \\[0.5em]
\hline
Fast & Slow \\[0.2em]
Pre-shared key & No pre-shared key \\[0.2em]
Owner(s) of the key & Owner of the private key \\[0.2em]
\end{tabular}
\end{large}
\end{center}
\note{
\begin{itemize}
\item \textbf{Symmetric}: faster, needs a pre-shared key
	\newline All people holding the key can decrypt the message.
\item \textbf{Asymmetric}: no need for a pre-shared key
	\newline Message can be decrypted by the person holding the right key.
\end{itemize}
\begin{block}{Usual usage}
A connection starts using an \emph{asymmetric} algorithm to exchange keys and
then switches to a \emph{symmetric} algorithm as it is faster.
\end{block}
}
\end{frame}

\begin{frame}
\frametitle{Signatures}
\includeimage[width=\linewidth]{signature-schema}
\note{
\begin{block}{Hash}
Fixed-size string built using an irreversible function from a message.
It's simple to use and a hash has a relatively small size. 
Signing the hash is more convenient than signing the whole message.
\newline Guarantees \emph{integrity}.
\end{block}
\begin{block}{Signature}
Mean to authenticate the identity of the sender of a message.
\end{block}
}
\end{frame}

%%%%%%%%%%%%%%%%%%%%%%%%%%%%%%%%%%%%%%%%%%%%%%%%%%%%%%%%%%%%%%%%%%%%%%%%%%%%%%

%------------------------------------------------
\subsection{Different means of authentication}
\begin{frame}
\begin{center}
\Huge{\textbf{Authentication means}}
\end{center}
\end{frame}

%------------------------------------------------
\begin{frame}
\frametitle{Password}
\includeimage[width=0.6\linewidth]{login}
\note{
\begin{itemize}
\item Password is the most common form of authentication
\item Password should be easily changeable
\item People should have different passwords for different services
\item Complexity must be sufficient
\end{itemize}
}
\end{frame}

\begin{frame}
\frametitle{Password complexity}
\includeimage[height=0.8\textheight,width=0.8\linewidth]{browserPasswordValidation}
\note{
Enforcing a minimum password strength:
\begin{itemize}
\item State the rules clearly (minimum 10 characters, at least one capital
letter, etc.)
\item Check the complexity in the browser and on the server
% feedback for the user
\end{itemize}
}
\end{frame}

\begin{frame}
\frametitle{Password strength}
\includeimage[height=0.75\textheight,width=\linewidth]{xkcd936_password_strength}
\note{
Explain the comic strip orally 
A passphrase vs a common password.
}
\end{frame}

%------------------------------------------------
\begin{frame}
\frametitle{One time passwords}
\includeimage[width=0.8\linewidth,height=0.7\textheight]{one-time-password}
\note{
\textbf{One time passwords} (OTP) are passwords which are valid only for a
single use.
\begin{itemize}
\item Prevent replay attack
\item Main difficulty: giving the users their passwords
	{\small (logistical problems)}
\item Used mainly by those who can afford big infrastructure
	\begin{itemize}
	\item states
	\item banks
	\item big companies
	\end{itemize}
\end{itemize}
}
\end{frame}

\begin{frame}
\frametitle{One time passwords - choice and distribution}
\begin{columns}
\begin{column}{0.5\linewidth}
\includeimage[width=\linewidth,height=0.7\textheight]{dice}
\end{column}
\begin{column}{0.5\linewidth}
\includeimage[width=\linewidth,height=0.7\textheight]{secure-id-token}
\end{column}
\end{columns}
\note{
To avoid prediction, passwords must be chosen using a random or at
least pseudo-random way.
\begin{itemize}
\item Time-based generation: passwords valid for a short time
\item Sent by SMS, e-mail, etc.
\item Software on a mobile phone
\end{itemize}
}
\end{frame}

%------------------------------------------------
\begin{frame}
\frametitle{Certificates}
\begin{columns}
\begin{column}{0.5\linewidth}
\includeimage[scale=0.5]{certificatesVariety1}
\end{column}
\begin{column}{0.5\linewidth}
\includeimage[scale=0.5]{certificatesVariety2}
\end{column}
\end{columns}
\includeimage[scale=0.5]{certificatesVariety3}
\note{
A trusted authority (Certification Authority) issues certificates to
confirm the identity of something. Those certificates may be of varying
quality and are most often used in SSL/TLS by web browsers.
}
\end{frame}

\begin{frame}
\frametitle{Certificate Authority}
\begin{itemize}
\item Sign, send and maintain certificates and revocation list
\end{itemize}
\end{frame}

\defverbatim\LstA{
\begin{lstlisting}
Certificate:
	Data:
		Version: 1 (0x0)
		Serial Number: 7829 (0x1e95)
		Signature Algorithm: md5WithRSAEncryption
		Issuer: C=ZA, ST=Western Cape,...
		        OU=Certification Services Division,
		        CN=Thawte Server CA/emailAddress=...
		Validity
			Not Before: Jul  9 16:04:02 1998 GMT
			Not After : Jul  9 16:04:02 1999 GMT
\end{lstlisting}
}
\defverbatim\LstB{
\begin{lstlisting}
		Subject: C=US, ST=Maryland,...
		         OU=FreeSoft, CN=www.freesoft.org/emailAddress=...
		Subject Public Key Info:
			Public Key Algorithm: rsaEncryption
			RSA Public Key: (1024 bit)
				Modulus (1024 bit):
					00:be ...:41:8f
				Exponent: 65537 (0x10001)
	Signature Algorithm: md5WithRSAEncryption
		93:5f...:68:9f
\end{lstlisting}
}
\begin{frame}
\frametitle{Certificates - What are they made of?}
\LstA
\end{frame}
\begin{frame}
\LstB
\end{frame}

\begin{frame}
\frametitle{Obtaining a certificate}
\includeimage[width=0.8\linewidth,height=0.75\textheight]{certification-process}
\note{
\begin{itemize}
\item When a website owner wants a certificate, he generates a public-private
key pair and then sends it to a CA.
\item The CA checks the integrity of the request and performs some checks on
the identity of the owner. Finally, it signs the certificate.
\end{itemize}
}
\end{frame}

\begin{frame}
\frametitle{Certificates - How are they used?}
\includeimage[scale=0.7]{one-way-ssl}
\note{
\begin{block}{One-way authentication}
The browser authenticates the server but the opposite isn't true.
\newline Common on the web.
\end{block}
}
\end{frame}

\begin{frame}
\frametitle{Certificates - How are they used?}
\includeimage[scale=0.7]{two-way-ssl}
\note{
\begin{block}{Two-way authentication}
Both ends are authenticated.
\end{block}
}

\note{
Certificates are often used as a mean to distribute keys in a public
key infrastructure (PKI).
A typical example is for SSL/TLS in web browsers.
\begin{itemize}
\item Each CA does generate it's public-private key pair. They
  self-sign their certificates with their public key.
\item Browsers are distributed with some pre-installed certificates.
  You can also add them manually.
\item When websites want to sign a message, they sign it with
  theirs private key and send the signed object and theirs certificates.
\item The browser checks the certificates against the CA
  certificate and the message against the checked certificate.
\end{itemize}
}
\end{frame}


\begin{frame}
%TODO : smaller image, Explain N + ONCE
\frametitle{Signature challenge}
\begin{columns}
\begin{column}{0.6\linewidth}
\includeimage[width=\linewidth,height=0.7\textheight]{challenge-accepted}
\end{column}
\begin{column}{0.4\linewidth}
\centering\Large\bfseries
nonce
\end{column}
\end{columns}
\note{
\begin{exampleblock}{A nonce}
Nonce stands for number used once. It's a random number used once to avoid
replay attacks.
\par
I encrypt a nonce with my correspondent's public key and ask him for the
decrypted and encrypted message.
\end{exampleblock}
This form of authentication is often used to sign online transactions.
}
\end{frame}

%------------------------------------------------
\begin{frame}
\frametitle{Tokens}
\includeimage[width=0.7\linewidth,height=0.6\textheight]{token}
\note{
\textbf{Token} are something of which ownership gives a form of
authentication.
\begin{exampleblock}{Forms in everyday life}
\begin{itemize}
\item Keys
\item Bank cards
\item Badges
\item Digipass
\end{itemize}
\end{exampleblock}
}
\end{frame}

%------------------------------------------------
\begin{frame}
\frametitle{Multi-factor authentication}
\begin{columns}[T]
\begin{column}{0.3\linewidth}
\begin{center}
\bfseries
IS
\includeimage[width=\linewidth]{multifactor-is}
\end{center}
\end{column}
\begin{column}{0.3\linewidth}
\begin{center}
\bfseries
OWNS
\includeimage[width=\linewidth]{multifactor-own}
\end{center}
\end{column}
\begin{column}{0.3\linewidth}
\begin{center}
\bfseries
KNOWS
\includeimage[width=\linewidth]{multifactor-know}
\end{center}
\end{column}
\end{columns}
\note{
\begin{center}
More than one form of authentication for an application.
\end{center}
The entity is authenticated by:
\begin{itemize}
\item What it \emph{is} (biometrics such as voice, fingerprints, ...)
\item What it \emph{owns} (a bank card, ...)
\item What it \emph{knows} (a password, a pin code, ...)
\end{itemize}
}
\end{frame}

%------------------------------------------------
\subsection{Single sign-on}

%------------------------------------------------
\begin{frame}
\frametitle{Single sign-on - principle}
\includeimage[width=0.8\linewidth,height=0.75\textheight]{sso}
\note{
\textbf{Single sign-on} (SSO) is a property of access control of multiple
related but independent software systems. With this property a user
logs in once and gains access to all systems without being prompted to
log in again for each of them.
\textit{-- Wikipedia}
\begin{columns}
\begin{column}{0.5\linewidth}
\begin{block}{Pro}
\begin{itemize}
\item Easier for the user
\item Implementations already exist
\item If they have only one password users tend to treat it with more care
\item Enter the password less often
\end{itemize}
\end{block}
\end{column}
\begin{column}{0.5\linewidth}
\begin{block}{Cons}
\begin{itemize}
\item All your eggs in the same basket (impact if compromised)
\item You are dependent upon your provider (confidence, availability, ...)
\item Loss of anonymity using SSO with public names
\end{itemize}
\end{block}
\end{column}
\end{columns}
}
\end{frame}

%------------------------------------------------
\subsection{Kerberos}

%------------------------------------------------
%TODO : too small
\begin{frame}
\frametitle{Kerberos}
\includeimage[width=3cm]{logo-kerberos}
\note{
\textbf{Kerberos} is a SSO developed at MIT to solve the problem of allowing
some users to use restricted resources.
MIT provides a free implementation of the protocol but it's also found in many
commercial products.
}
\end{frame}

\begin{frame}
%TODO : too small and not clear
\frametitle{Kerberos - How does it work?}
\includeimage[width=0.8\linewidth]{kerberos}
\note{
\begin{enumerate}
\item The client gets a \emph{ticket} for a distribution center to
authenticate itself
\item The client gets a \emph{service ticket}
\item The \emph{service ticket} can be used on another server
\end{enumerate}
}
\end{frame}

%------------------------------------------------
\subsection{Provider extensions}

%------------------------------------------------
\begin{frame}
%TODO : move to before Kerberos
\frametitle{Provider extensions}
\begin{columns}
\begin{column}[width=\linewidth]{0.3\linewidth}
\includeimage[width=0.8\linewidth]{openid}
\end{column}
\begin{column}[width=\linewidth]{0.3\linewidth}
\includeimage[width=0.8\linewidth]{logo-google}
\end{column}
\begin{column}[width=\linewidth]{0.3\linewidth}
\includeimage[width=0.8\linewidth]{logo-facebook}
\end{column}
\end{columns}
\note{
\begin{itemize}
\item Single sign-on provided using:
	\begin{itemize}
	\item Facebook account
	\item Google
	\item OpenID
		\begin{itemize}
		\item myOpenID
		\item ClaimID
		\item ...
		\end{itemize}
	\end{itemize}
\item Simpler to setup than developing your own single sign-on mechanism
\end{itemize}
}
\end{frame}

%------------------------------------------------
\subsection{Compromised authentication mean}

%------------------------------------------------
%TODO : image too dark
\begin{frame}
\frametitle{Compromised authentication mean}
\includeimage[width=0.8\linewidth]{lost-credentials}
\note{
When you implement authentication you have to take into account that
the user may lose its authentication mean or it can get stolen.
\begin{itemize}
\item Deactivate a compromised authentication mean
\item Authenticate the user by an uncompromised mean
\item Give him a new primary authentication mean
\end{itemize}
}
\end{frame}

%------------------------------------------------
\begin{frame}
%TODO : too small and add side channels (post, sms,...)
\frametitle{Example: Lost password}
\begin{itemize}
\item Security questions
\item Another channel {\small (e-mail)}
\end{itemize}
\note{
If one of your users loses his password, he has to get a new one but doing
this in a secure way can be challenging.
\begin{itemize}
\item Use some predefined security questions
\item Send a token over a side channel {\small (e-mail)}
\item Confirm the change
\end{itemize}
\begin{block}{Challenges}
An attacker may use social engineering to get answers to security questions.
\end{block}
}
\end{frame}



\section{Attacks}

%%%%%%%%%%%%%%%%%%%%%%%%%%%%%%%%%%%%%%%%%%%%%%%%%%%%%%%%%%%%%%%%%%%%%%%%%%%%%%
\subsection{Attack trees}

\begin{frame}
\includeimage[width=\linewidth,height=6cm]{attack-tree-admin}
\note{
\begin{block}{What is it?}
Model describing how an \emph{asset} can be attacked.
\end{block}
\begin{block}{Goal}
Used to determine and understand threats that may arise.
\end{block}
}
\end{frame}

%%%%%%%%%%%%%%%%%%%%%%%%%%%%%%%%%%%%%%%%%%%%%%%%%%%%%%%%%%%%%%%%%%%%%%%%%%%%%%
\subsection{Threat agents}

\begin{frame}
\begin{columns}
\begin{column}{0.6\linewidth}
\includeimage[width=\linewidth,height=5cm]{hacker}
\end{column}
\begin{column}{0.4\linewidth}
\includeimage[width=\linewidth,height=3cm]{yes-i-can}
\end{column}
\end{columns}
\note{
\begin{description}
\item[External threats]
	\hfill
	\begin{itemize}
	\item internet users
	\item viruses
	\end{itemize}
\item[Internal threats]
	Unsatisfied employees (intranet users)
\item[Reason]
	To show off, break things
\end{description}
}
\end{frame}

\begin{frame}
\begin{columns}
\begin{column}{0.6\linewidth}
\includeimage[width=\linewidth]{organized-crime}
\end{column}
\begin{column}{0.4\linewidth}
\includeimage[width=0.6\linewidth]{money}
\end{column}
\end{columns}
\note{
\begin{description}
\item[External threats]
	\hfill
	\begin{itemize}
	\item organized crime
	\item competitor
	\end{itemize}
\item[Reason]
	Gain
	\begin{itemize}
	\item financial
	\item information
	\item competition
	\end{itemize}
\end{description}
\begin{block}{Impact}
Threat agents have different skills, resources and different motivations that
may influence risk.
\end{block}
}
\end{frame}

\begin{frame}
\begin{columns}
\begin{column}{0.5\linewidth}
\includeimage[width=\linewidth,height=4cm]{crew}
\end{column}
\begin{column}{0.5\linewidth}
\includeimage[width=\linewidth]{disaster}
\end{column}
\end{columns}
\note{
\begin{description}
\item[Internal threats]
	Employees, crew
\item[Natural threats]
	(disasters)
\item[Technical failure]
	\hfill
	\begin{itemize}
	\item loss of essential services (electricity, etc.)
	\item hardware failure
	\item software failure
	\end{itemize}
\item[Reason]
	Unintentional, accident
\end{description}
}
\end{frame}

%%%%%%%%%%%%%%%%%%%%%%%%%%%%%%%%%%%%%%%%%%%%%%%%%%%%%%%%%%%%%%%%%%%%%%%%%%%%%%
\subsection{Risk}

\begin{frame}
\frametitle{What is a \emph{risk}?}
\begin{quote}
The potential that a given threat will exploit vulnerabilities of
an asset or group of assets and thereby cause harm to the organization. It
is measured in terms of a combination of the probability of occurrence of an
event and its consequence.
\textit{-- Wikipedia}
\end{quote}
\note{
Just a good definition.
Explanation in the next slide.
}
\end{frame}

\begin{frame}
\begin{Huge}
\[ \text{risk} = \text{likelihood} * \text{impact} \]
\end{Huge}
\note{
\begin{itemize}
\item Likelihood: probability of a successful attack
	\\ Influenced by threat agent skills and vulnerability factors
\item Impact: how much damage the attack causes
\end{itemize}
\begin{columns}
\begin{column}{0.5\linewidth}
\par\textbf{Vulnerability factors}
\begin{itemize}
\item Ease of discovery
\item Ease of exploit
\item Awareness
\item Intrusion detection
\end{itemize}
\end{column}
\begin{column}{0.5\linewidth}
\par\textbf{Impact factors}
\begin{itemize}
\item Loss of confidentiality, integrity, availability, accountability
\item Financial damage
\item Reputation damage
\item Privacy violation
\end{itemize}
\end{column}
\end{columns}
}
\end{frame}

%%%%%%%%%%%%%%%%%%%%%%%%%%%%%%%%%%%%%%%%%%%%%%%%%%%%%%%%%%%%%%%%%%%%%%%%%%%%%%
\subsection{Context of web applications}

\subsubsection{Network}

\begin{frame}
\includeimage[width=0.8\linewidth]{client-server}
\note{
\begin{itemize}
\item Applications available on the internet
\item Anyone can access the application
\end{itemize}
\begin{block}{Anonymity}
No identifier for a web user.
\\ Example: IP address shared by multiple users
\par\textbf{Session number}\\
The server identifies its users using a session ID (unique number) for each
client.
\end{block}
}
\end{frame}

\subsubsection{Architecture of a web application}

\begin{frame}
\includeimage[width=\linewidth,height=0.8\textheight]{archi-webapplication}
\note{
\begin{itemize}
\item Browser: the client, uses scripting
\item Router: connected to the browser (by WiFi) - LAN
\item Internet: not detailed here - WAN
\item DMZ: demilitarized zone
\item Web Application Firewall, IDS, Load Balancer
\item Application Server: with the contained application
\item Database: where the application stores its data
\end{itemize}
\begin{block}{Points of failure}
Any component may fail; will the web site continue to work?
\begin{description}
\item[Single point of failure]
A part of a system that prevents the entire system from working when it fails.
Solution: duplication.
\end{description}
\end{block}
}
\end{frame}

%%%%%%%%%%%%%%%%%%%%%%%%%%%%%%%%%%%%%%%%%%%%%%%%%%%%%%%%%%%%%%%%%%%%%%%%%%%%%%
\subsection{OWASP - Top 10}

\begin{frame}
\centering
{\bfseries\usebeamercolor[fg]{block title}O}pen
{\bfseries\usebeamercolor[fg]{block title}W}eb
{\bfseries\usebeamercolor[fg]{block title}A}pplication
{\bfseries\usebeamercolor[fg]{block title}S}ecurity
{\bfseries\usebeamercolor[fg]{block title}P}roject
\includeimage[width=0.5\linewidth,height=3cm]{logo-owasp}
\note{
\begin{block}{OWASP}
\begin{itemize}
\item Standards
\item Libraries
\item Books
\item Etc.
\end{itemize}
\end{block}
\begin{block}{Project: Top 10}
OWASP compiles a top 10 of most critical web application \emph{risks}.
\end{block}
}
\end{frame}

\begin{frame}
\frametitle{Top 10}
\begin{enumerate}
\item Injection
\item Cross-site Scripting (XSS)
\item Authentication and Session Management
\item Insecure Direct Object References
\item Cross-site Request Forgery (CSRF)
\item Security Misconfiguration
\item Insecure Cryptographic Storage
\item Failure to Restrict URL Access
\item Insufficient Transport Layer Protection
\item Unvalidated Forwards and Redirects
\end{enumerate}
\note{
\begin{block}{Project: Top 10}
Top 10 of most critical web application \emph{risks}.
\end{block}
\begin{block}{Demo application}
\begin{itemize}
\item Built with security flaws for the demo
\item Hotel rating web site
\item Shown hotels must be approved by an administrator
\end{itemize}
\end{block}
}
\end{frame}

\subsubsection{Attack 1 - Injection}

\begin{frame}
\frametitle{1. Injection}
%\includeimage[width=0.8\linewidth,height=0.75\textheight]{syringe}
\includeimage[width=0.8\linewidth]{client-server-injection}
\note{
\textbf{Injection} consists in sending untrusted data to an interpreter.
Attacker enters code to execute into input fields.
\begin{example}
Send code directly to some interpreter.
\end{example}
\begin{block}{Impacts}
Data loss, corruption, lack of accountability, etc.
\end{block}
\begin{alertblock}{Risks}
Easy to exploit, has a severe impact, it is quite common.
\end{alertblock}
}
\end{frame}

\defverbatim[colored]\Lst{
\begin{lstlisting}[style=beamer]
String query = "select * from users"
	+ " where user_name = '" @+ name +@ "'"
	+ " and password = '" @+ password +@ "'";
\end{lstlisting}
}
\begin{frame}
\begin{exampleblock}{\emph{THE} password}
\includeimage[height=0.4\textheight]{the-password}
\pause
\begin{embeddedcode}{bad}
\Lst
\end{embeddedcode}
\end{exampleblock}
\note{
\begin{exampleblock}{SQL injection - demo}
Getting the users' passwords as administrator.
\end{exampleblock}
\begin{exampleblock}{JPQL injection - demo}
Searching hotels that have a manager that has the same password as a given
user.
\end{exampleblock}
}
\end{frame}

\defverbatim[colored]\Lst{
\begin{lstlisting}[style=beamer]
String query = "select * from users where @user_name = ?@";
PreparedStatement st = con.prepareStatement(query);
@st.setString(1, name);@
ResultSet rs = st.executeQuery();
\end{lstlisting}
}
\begin{frame}
\begin{embeddedcode}{good}
\Lst
\end{embeddedcode}
\pause
\begin{itemize}
\item Parametrized interface
\item Escaping routines
\item White list validation
\end{itemize}
\note{
\begin{block}{Using parametrized interface}
Usage of placeholders for values let's the interpreter escape input values.
\end{block}
\begin{block}{Escaping routines}
Some languages have special functions to escape manually values.
\end{block}
\begin{block}{White list validation}
List of valid input patterns.
\end{block}
}
\end{frame}

\subsubsection{Attack 2 - XSS}

\begin{frame}
\frametitle{2. XSS}
\includeimage[width=0.8\linewidth]{client-server-xss}
\note{
\textbf{Cross-site scripting} allows attackers to inject code into the page
sent to the user. The injected code can be a script interpreted by the
browser.
\begin{example}
The attacker sends a script to a web page. Users who download the code execute
it in their browser.
\end{example}
\begin{block}{Impacts}
Hijack user sessions, change content, redirect the user.
\end{block}
\begin{alertblock}{Risks}
The most widespread vulnerability.
It requires an average knowledge to be exploited and the impacts are moderate.
\end{alertblock}
}
\end{frame}

\defverbatim[colored]\LstBad{
\begin{lstlisting}[style=beamer]
<div>
	@${hotel.description}@
</div>
\end{lstlisting}
}
\defverbatim[colored]\LstGood{
\begin{lstlisting}[style=beamer]
<div>
	@<c:out value="${hotel.description}" />@
</div>
\end{lstlisting}
}
\begin{frame}
\begin{embeddedcode}{bad}
\LstBad
\end{embeddedcode}
\badgoodsep
\begin{embeddedcode}{good}
\LstGood
\end{embeddedcode}
\note{
Unescaped included text in a JSP.
\begin{exampleblock}{Getting the user's cookies - demo}
Just use \lstinline!document.cookie! in JavaScript and send it somewhere.
\newline Explain what are \emph{cookies} orally.
\end{exampleblock}
\begin{block}{Escaping the values}
All the values should be escaped before sending them to the users.
\end{block}
\begin{block}{White list validation}
Output values can be white listed, but this is not a complete defence against
XSS as sometimes special characters must be accepted.
\end{block}
}
\end{frame}

\subsubsection{Attack 3 - Authentication and Session Management}

\begin{frame}
\frametitle{3. Authentication and Session Management}
\includeimage[width=0.8\linewidth]{session-hijacking}
\note{
The attacker uses flaws in the \textbf{authentication or session management}
implemented for a given web site to steal someone else's identity.
\begin{example}
\begin{itemize}
\item Session hijacking
\item Session fixation
\end{itemize}
\end{example}
\begin{block}{Impacts}
Once an account has been stolen, the attacker may do \emph{anything} the user
can do.
\end{block}
\begin{alertblock}{Risks}
The impact is severe, this attack is common.
\end{alertblock}
}
\end{frame}

\defverbatim[colored]\Lst{
\begin{lstlisting}[style=beamer]
http://example.com/page@;jsessionid=dfd4f5fda35df@...
\end{lstlisting}
}
\begin{frame}
\begin{columns}
\begin{column}{0.7\linewidth}
\includeimage[width=\linewidth,height=3cm]{cybercafe}
\end{column}
\begin{column}{0.3\linewidth}
\centering\Huge\bfseries
XSS
\end{column}
\end{columns}
\vspace{6mm}
\Lst
\badgoodsep
\includeimage[height=3cm]{irisscan}
\note{
\begin{exampleblock}{Public computer}
Forgetting to log out...
\end{exampleblock}
\begin{exampleblock}{Session fixation}
In older application server, the session ID could be set using the URL.
You could send that URL to someone else.
\end{exampleblock}
\pause
\begin{exampleblock}{XSS session hijacking - demo}
The user's session can be obtained using XSS.
\end{exampleblock}
\hrule
\begin{block}{Good authentication mechanism}
\begin{itemize}
\item Use a proven authentication mechanism
\item Prevent \emph{XSS attacks}
\item Good session timeouts and accessible log out buttons
\end{itemize}
\end{block}
}
\end{frame}

\subsubsection{Attack 4 - Insecure Direct Object References}

\defverbatim[colored]\Lst{
\begin{lstlisting}[style=beamer]
http://example.com/user/@152@
\end{lstlisting}
}
\begin{frame}
\frametitle{4. Insecure Direct Object References}
\begin{center}
\Lst
\end{center}
\note{
The attacker changes a parameter in the request to obtain a \textbf{direct
object reference} that they should not be able to access.
\begin{example}
Replace some object identifier in the request.
\end{example}
\begin{block}{Impacts}
Compromise the data that can be referenced.
\end{block}
\begin{alertblock}{Risks}
Easy to exploit and detect. The impact is moderate.
\end{alertblock}
}
\end{frame}

\begin{frame}
\begin{center}
{\Huge\bfseries Check permissions}
\end{center}
\note{
\begin{block}{Missing access verification}
Often objects are retrieved by their ID, the developer can forget to check
whether the user is allowed to see that object.
\end{block}
\hrule
\begin{exampleblock}{An unvalidated hotel - demo}
See an unvalidated hotel as a normal user.
\end{exampleblock}
\hrule
\begin{block}{\textbf{Check permissions}}
Check whether the user has the permissions to access the object before
manipulating it.
\end{block}
\begin{block}{References by session}
Use indirect references for each session.
\end{block}
}
\end{frame}

\subsubsection{Attack 5 - CSRF}

\begin{frame}
\frametitle{5. CSRF}
\includeimage[width=0.8\linewidth]{client-server-csrf}
\note{
\textbf{Cross-site Request Forgery} consists in generating requests when the
user visits the attacker's web page.
\begin{example}
\begin{itemize}
\item A user visits the attacker's web page
\item The requested page contains tricks to submit requests to the attacked
web site as the visiting user
\end{itemize}
\end{example}
\begin{block}{Impacts}
The attacker may perform actions as the victim.
\end{block}
\begin{alertblock}{Risks}
Widespread and easy to detect, the impact is moderate.
\end{alertblock}
}
\end{frame}

\defverbatim[colored]\LstBad{
\begin{lstlisting}[style=beamer]
<img src="http://@example.com@/forum/thread1/
@addComment@?text=hello%20world" alt="" />
\end{lstlisting}
}
\defverbatim[colored]\LstGood{
\begin{lstlisting}[style=beamer]
<form @method="post"@ action="addComment">
	<input @type="hidden"@ name="token" @value="<random>"@ />
	...
</form>
\end{lstlisting}
}
\begin{frame}
\begin{embeddedcode}{bad}
\LstBad
\end{embeddedcode}
\badgoodsep
\begin{embeddedcode}{good}
\LstGood
\end{embeddedcode}
\note{
\begin{exampleblock}{Post a comment}
The attacker's page may contain:
\end{exampleblock}
\hrule
\begin{block}{Unpredictable token}
Each request contains a field that is a token linked to the session of
the user, the attacker's website will not be able to generate the correct
requests.
\newline
Tokens can be unique by session or by request.
\par
The token is usually stored in a cookie (preferably HTTP-only).
In forms, a hidden field contains the same value.
When the request is sent, both values are compared.
\par
Only the session ID is stored in the cookie; the session information is
stored only on the application server.
\end{block}
}
\end{frame}

\subsubsection{Attack 6 - Security Misconfiguration}

\begin{frame}
\frametitle{6. Security Misconfiguration}
\includeimage[width=0.7\linewidth]{misconfiguration}
\note{
\textbf{Security configuration} can happen at any level: the application, the
web server, the framework, etc.
The default behavior is not always the wanted one.
\begin{example}
\begin{itemize}
\item Accessing: default accounts, unprotected files and directories
\item Gaining information: versions, stack traces, debug messages
\end{itemize}
\end{example}
\begin{block}{Impacts}
Attackers can gain unauthorized access, the system can be compromised.
\end{block}
\begin{alertblock}{Risks}
Easy to exploit and to detect.
\end{alertblock}
}
\end{frame}

\begin{frame}
\includeimage[height=0.7\textheight]{box}
\note{
The out-of-the-box configuration is not always the best.
}
\end{frame}

\begin{frame}
\includeimage[width=\linewidth,height=0.7\textheight]{webpage-stacktrace}
\pause
\raggedleft\Large
... more information for the attacker
\note{
\begin{block}{Outdated libraries}
Bugs are published for known libraries and applications.
\end{block}
\begin{exampleblock}{Stack trace}
A stack trace may contain the used parameters, references to libraries, etc.
\end{exampleblock}
\hrule
\begin{block}{Up-to-date system}
Keep the used software up to date and manage versions.
\end{block}
\begin{block}{Good architecture}
Good separation between components is more secure.
\end{block}
\begin{block}{Automated scan}
Automated scans can detect misconfigurations and missing patches.
\end{block}
}
\end{frame}

\subsubsection{Attack 7 - Insecure Cryptographic Storage}

\begin{frame}
\frametitle{7. Insecure Cryptographic Storage}
\begin{center}
\begin{tabular}{|l|l|l|c|}
\hline
ID & NAME & PASSWORD & \ldots \\
\hline
1 & admin & 2ezf5zf2d & \ldots \\
2 & jacky & toto & \ldots \\
3 & k & 123456 & \ldots \\
\hline
\end{tabular}
\end{center}
\note{
Sensitive data is stored \textbf{insecurely without crypting}.
\begin{example}
All data is readable by the attacker once he gains access to it.
\end{example}
\begin{block}{Impacts}
Stolen data is compromised.
\end{block}
\begin{alertblock}{Risks}
All unencrypted data is compromised.
The attack is however difficult to perform.
\end{alertblock}
}
\end{frame}

\begin{frame}
\includeimage[height=0.6\textheight,width=0.8\linewidth]{rainbow-table-example}
\note{
\begin{exampleblock}{Unencrypted passwords}
After dumping the table \texttt{USERS}, the passwords might be stored in
plain text.
\end{exampleblock}
\begin{exampleblock}{Rainbow tables}
Precomputed table for decoding hashed values.
\href{http://www.md5this.com/list.php?page=133&key=1}{\texttt{md5this.com}}
\end{exampleblock}
}
\end{frame}

\begin{frame}
\includeimage[height=0.6\textheight,width=0.8\linewidth]{backup}
\note{
\begin{exampleblock}{Unencrypted backup}
The database backup is lost.
\end{exampleblock}
}
\end{frame}

\begin{frame}
\frametitle{Hash + Salt}
\includeimage[height=0.7\textheight,width=\linewidth]{hash-salt}
\note{
Ensure passwords are hashed with a standard algorithm and a different salt is
used for each password.
\par
Prevents rainbow table attacks.
\par
Both the hash and salt are stored in a database.
}
\end{frame}

\subsubsection{Attack 8 - Failure to Restrict URL Access}

\defverbatim[colored]\Lst{
\begin{lstlisting}[style=beamer]
http://example.com/@admin@
\end{lstlisting}
}
\begin{frame}
\frametitle{8. Failure to Restrict URL Access}
\Lst
\note{
\textbf{URL access is not restricted} to functions requiring permissions.
\begin{example}
The attacker changes his URL to a privileged page.
\end{example}
\begin{block}{Impacts}
Administrative functions may become available to the attacker.
\end{block}
\begin{alertblock}{Risks}
Easy to exploit, but the flaw is uncommon.
\end{alertblock}
}
\end{frame}

\begin{frame}
\begin{center}
{\Huge\bfseries Authorization}
\end{center}
\note{
\begin{exampleblock}{Administration page}
\url{http://example.com/admin} is not referenced from any other page, but
the authentication is not enforced.
\end{exampleblock}
\hrule
\begin{block}{Check permissions}
Authorization mechanisms should be used.
\end{block}
}
\end{frame}

\subsubsection{Attack 9 - Insufficient Transport Layer Protection}

\begin{frame}
\frametitle{9. Insufficient Transport Layer Protection}
\includeimage[width=0.8\linewidth]{wifi-sniffing}
\note{
\textbf{Traffic is unprotected} and anyone sniffing the network can see the
data in transit.
\begin{example}
\begin{itemize}
\item Sniffing the WiFi network
\item Scanning for HTTP requests
\item Analyzing conversations
\end{itemize}
\end{example}
\begin{block}{Impacts}
Account theft, expose individual user's data.
\end{block}
\begin{alertblock}{Risks}
Detectable using any sniffer.
However, often difficult to exploit.
\end{alertblock}
}
\end{frame}

\begin{frame}
\includeimage[width=0.5\linewidth]{https}
\note{
The attacker can see the data that is sent over the network.
\begin{exampleblock}{Sniffing passwords}
Sniff the network and find HTTP requests containing "password".
\end{exampleblock}
\hrule
\begin{block}{HTTP\textbf{S}}
SSL/TLS can be used to secure the connection.
\begin{itemize}
\item \emph{secure} cookies will not be sent through an unencrypted connection
\item Ensure that certificates are valid
\end{itemize}
\end{block}
}
\end{frame}

\subsubsection{Attack 10 - Unvalidated Forwards and Redirects}

\defverbatim[colored]\Lst{
\begin{lstlisting}[style=beamer]
http://example.com/forward.jsp?@fwd=admin.jsp@
\end{lstlisting}
}
\begin{frame}
\frametitle{10. Unvalidated \textbf{Forwards} and Redirects}
\includeimage[scale=0.5]{webapp-forward}
\vspace{6mm}
\Lst
\note{
A URL may appear as trustworthy, but it contains an \textbf{unvalidated
forward or redirect}.
\begin{example}
The attacker may bypass authorization for \texttt{admin.jsp}.
\end{example}
\begin{block}{Impacts}
Tricking victims into disclosing information, unsafe forwards may allow access
control bypass.
\end{block}
\begin{alertblock}{Risks}
Easy to detect, uncommonly present.
\end{alertblock}
}
\end{frame}

\defverbatim[colored]\Lst{
\begin{lstlisting}[style=beamer]
http://@bank.com/secured@/../redirect
?@url=http://evil-bank.com/secured@
\end{lstlisting}
}
\begin{frame}
\frametitle{10. Unvalidated Forwards and \textbf{Redirects}}
\includeimage[scale=0.5]{webapp-redirect}
\vspace{6mm}
\Lst
\note{
\begin{example}
Placing a URL with a redirect that seems trustworthy redirecting to a
malicious page.
\end{example}
\begin{exampleblock}{Malicious site}
Redirect to a malicious website using:
\\ \url{http://bank.com/redirect?url=evil-bank.com}
\end{exampleblock}
\hrule
\begin{block}{Avoid redirects and forwards}
Avoid \emph{parametrized} redirects and forwards.
\end{block}
\begin{block}{Authorization}
Validate and authorize the user after each redirect.
\end{block}
\begin{exampleblock}{Map input values}
Map input values to the actual URLs.
\end{exampleblock}
}
\end{frame}

%%%%%%%%%%%%%%%%%%%%%%%%%%%%%%%%%%%%%%%%%%%%%%%%%%%%%%%%%%%%%%%%%%%%%%%%%%%%%%
\subsection{DoS}

\begin{frame}
\frametitle{Denial of Service (DoS)}
\includeimage[width=0.8\linewidth,height=0.8\textheight]{overwhelmed-donkey}
\note{
Attack that aims to bring down (or make unavailable) some service usually by
causing server overload.
\begin{itemize}
\item SYN flood
\item Application level flood
\item R-U-Dead-Yet
	\newline
	(never-ending POST transmissions,
	example: \href{http://ha.ckers.org/slowloris/}{Slowloris})
\item DDoS
\item Unintentional DoS
\end{itemize}
}
\end{frame}

\begin{frame}
\begin{columns}
\begin{column}{0.5\linewidth}
\includeimage[width=\linewidth,height=0.5\textheight]{anonymous-group}
\end{column}
\begin{column}{0.5\linewidth}
\includeimage[width=\linewidth]{dos-unintentional}
\end{column}
\end{columns}
\note{
\frametitle{Denial of Service - example}
\begin{exampleblock}{DDoS}
The anonymous group attack.
\end{exampleblock}
\begin{exampleblock}{Unintentional DoS}
\begin{itemize}
\item page linked from \url{http://slashdot.org}
\item \href{http://www.joueurdugrenier.fr/}{Joueur du Grenier} - new video
\end{itemize}
\end{exampleblock}
\hrule
\begin{block}{Denial of Service - prevention}
\begin{itemize}
\item Firewall
\item Intrusion Prevention System
\item Black-holing
\end{itemize}
\end{block}
}
\end{frame}

%%%%%%%%%%%%%%%%%%%%%%%%%%%%%%%%%%%%%%%%%%%%%%%%%%%%%%%%%%%%%%%%%%%%%%%%%%%%%%
\subsection{MitM}

\begin{frame}
\frametitle{Man in the Middle attack}
\includeimage[height=0.6\textheight]{phone-tapping}
\note{
The attacker plays the role of the server to the client and the role of the
client to the server.
All communication transits through the attacker who can modify messages as
they are sent.
\begin{exampleblock}{Public WiFi}
Vulnerable on public WiFi networks.
\\ Tools such as \texttt{ettercap} allow to perform the attack.
\end{exampleblock}
\hrule
\begin{block}{Prevention}
Use cryptography: HTTPS, PKI (Public Key Infrastructure)
\end{block}
}
\end{frame}

%%%%%%%%%%%%%%%%%%%%%%%%%%%%%%%%%%%%%%%%%%%%%%%%%%%%%%%%%%%%%%%%%%%%%%%%%%%%%%
\subsection{Social engineering}

\begin{frame}
\frametitle{Social engineering}
%\includeimage[width=0.8\linewidth]{social-engineering}
\includeimage[width=0.8\linewidth]{networking-sites}
\note{
\textbf{Social engineering} consists in manipulating people so they divulge
confidential information.
\\ Similar to fraud.
}
\end{frame}

\begin{frame}
\includeimage[width=0.8\linewidth,height=0.6\textheight]{hannibal-lecter}
\begin{quote}
\begin{center}
Quid pro quo, Clarice...
\end{center}
\end{quote}
\note{
\begin{description}
\item[Pretexting]
Using an invented scenario (the pretext) to engage a targeted victim in
a manner that increases the chance the victim will divulge information
or perform actions that would be unlikely in ordinary circumstances.
\item[Phishing]
See the next point.
\item[Quid pro quo]
Something for something:
Help a user solve a problem, they might give you passwords in exchange.
\end{description}
}
\end{frame}

\begin{frame}
\begin{columns}
\begin{column}{0.4\linewidth}
\Large
\centering
\emph{password}
\end{column}
\begin{column}{0.2\linewidth}
\includeimage[height=8mm]{vs}
\end{column}
\begin{column}{0.4\linewidth}
\includeimage[width=\linewidth,height=3cm]{pen}
\end{column}
\pause
\end{columns}
\begin{exampleblock}{Infosecurity Europe 2003 survey}
\begin{quote}
In a 2003 information security survey, {\Large 90\%} of office workers gave
researchers what they claimed was their {\Large password} in answer to a
survey question in exchange for a {\Large cheap pen}.
\end{quote}
\end{exampleblock}
\note{
\begin{block}{Protection}
\begin{itemize}
\item Education
\item Framework of trust
\end{itemize}
\end{block}
}
\end{frame}

%%%%%%%%%%%%%%%%%%%%%%%%%%%%%%%%%%%%%%%%%%%%%%%%%%%%%%%%%%%%%%%%%%%%%%%%%%%%%%
\subsubsection{Phishing}

\begin{frame}
\frametitle{Phishing}
\includeimage[height=0.75\textheight,width=\linewidth]{phishing-mail}
\note{
\textbf{Phishing} is the attempt to acquire user information through
masquerading as a legitimate web site.
\begin{itemize}
\item Cloned websites
\item Link manipulation
\item Phone phishing
\item Etc.
\end{itemize}
\hrule
\begin{block}{Phishing - prevention}
\begin{itemize}
\item Education
\item Help to identify legitimate web sites
\end{itemize}
\end{block}
}
\end{frame}

\begin{frame}
\includeimage[height=0.75\textheight,width=\linewidth]{cloned-website}
\note{
A \emph{cloned} website.
}
\end{frame}

%%%%%%%%%%%%%%%%%%%%%%%%%%%%%%%%%%%%%%%%%%%%%%%%%%%%%%%%%%%%%%%%%%%%%%%%%%%%%%
\subsection{Other attacks}

\begin{frame}
\frametitle{Brute force}
%\includeimage[height=0.7\textheight]{letstryallthepasswords}
\includeimage[height=0.7\textheight]{hulk-smash}
\note{
\begin{block}{Try all the passwords}
Try all the possible combinations.
\newline
$k^n$ where $k$ is the size of the alphabet and $n$ the length of the
password.
\end{block}
\begin{block}{Dictionary attack}
Try a list of common passwords.
\\ English language contains only hundreds of thousands words.
\end{block}
}
\end{frame}

\defverbatim[colored]\Lst{
\begin{lstlisting}[style=beamer]
http://example.com/forum/download.jsp
?@path=/etc/shadow@
\end{lstlisting}
}
\begin{frame}
\frametitle{Path traversal}
\Lst
\note{
An input is a relative unvalidated path.
You can download or upload any file that the server contains.
\begin{exampleblock}{Download form}
Get a file uploaded to a forum... or the password file.
\end{exampleblock}
}
\end{frame}

\begin{frame}
\frametitle{Failure to restrict automation}
\begin{itemize}
\item Spamming
\item REST API
\end{itemize}
\pause
\vspace{2cm}
\includeimage[width=0.6\linewidth]{recaptcha-example}
\note{
The web site is not protected against automation.
\begin{itemize}
\item Common for REST API
\item Programs may overload the website unintentionally
	\newline
	(Fix: maximum number of requests per IP for a time period)
\end{itemize}
\begin{exampleblock}{Captcha - prevention}
Challenge-response test easy for humans, difficult for machines.
\end{exampleblock}
\begin{block}{Captcha - cracking}
An evil website could present a captcha to crack another website's protection.
Users would crack unknowingly some other website's protection against
automated attacks.
\end{block}
}
\end{frame}

\begin{frame}
\frametitle{And so on...}
\includeimage[width=0.8\linewidth,height=0.75\textheight]{attack-tag-cloud}
\note{
As technology evolves, new forms of attacks appear.
\par
Instead of trying to fix vulnerabilities, focus on establishing strong
security controls.
}
\end{frame}

%%%%%%%%%%%%%%%%%%%%%%%%%%%%%%%%%%%%%%%%%%%%%%%%%%%%%%%%%%%%%%%%%%%%%%%%%%%%%%
\subsection{Common points}

\begin{frame}
\frametitle{Common points}
\begin{itemize}
\item<+-> Improper input validation
\item<+-> Missing authorization checks
\item<+-> Misconfiguration
\item<+-> Other assumptions
\item<+-> Education factor {\small (social engineering)}
\end{itemize}
\note{
Bullet points.
}
\end{frame}



\section{Prevention}

%------------------------------------------------
\subsection{Different moments}

%------------------------------------------------
\begin{frame}
\frametitle{Different moments}

\begin{itemize}
\item Architecture
\item Development
\item Maintenance
\end{itemize}
\note{
Preventing security flaws is something that has to be done during the
whole life cycle of the application.
}
\end{frame}

%------------------------------------------------
\begin{frame}
\frametitle{Architecture}
% TODO: ajouter une image du livre?
\begin{quote}
\begin{small}
The {\Large highest-level breakdown} of a system into its parts; the decisions that are
{\Large hard to change}; there are multiple architectures in a system; what is
architecturally significant can change over a system's lifetime; and, in the
end, architecture boils down to whatever the {\Large important stuff} is.
\newline \textit{-- Martin Fowler, Patterns of Enterprise Application Architecture}
\end{small}
\end{quote}
\note{
Definition of \emph{architecture}.
}
\end{frame}

\begin{frame}
\frametitle{Security design\ldots}
\includeimage[width=0.7\linewidth,height=0.7\textheight]{useless-barrier}
\note{
\begin{itemize}
\item Security has to be taken into account from the beginning.
	\newline
	Retrofitting security in an unsecure application is much more difficult
	than conceiving it securely from the start.
\item You have to ask yourself the good questions and those questions vary
	from one application to the other
\item Know the threats you are facing
\item There is no miracle receipt... just guidelines
\end{itemize}
Question examples:
\begin{itemize}
\item Who are the end-users?
\item In what environment will the application run?
\item Who has access to the data?
\item Are there other systems that access data? What are the rules?
\end{itemize}
% Source:
% https://www.owasp.org/index.php/Application_Security_Architecture_Cheat_Sheet}
}
\end{frame}

%------------------------------------------------
\begin{frame}
\frametitle{Development: tips and tricks}
\begin{itemize}
\item<+-> Passwords
\item<+-> Session management
\item<+-> Input validation
\item<+-> Ad hoc protection
\end{itemize}
\note{
There are lots of best practices that can make your application more
secure.
\begin{itemize}
\item Passwords: enforcing a minimum complexity policy, storing them securely,
having the user rotate them, preventing brute-force by implementing temporary
lock-out, ...
\item Session management: setting an inactivity time out, setting the
\emph{secure} flag for sensitive cookies, implementing an easily accessible logout, ...
\item Input validation: always on server side, use a positive
  approach
\item Use ad-hoc protection for known attacks. Ad-hoc protection means a
protection specialized for a particular attack in this context.
\end{itemize}
}
\end{frame}

%------------------------------------------------
\begin{frame}
%TODO : changer le logo
\frametitle{Secure development cycle}
\includeimage[width=0.7\linewidth,height=0.7\textheight]{development-cycle}
\note{
A secure development cycle should set aside some time for:
\begin{itemize}
\item Code review {\small (security oriented and general)}
\item Pentesting
\item Security training of the developers
\item Sensitizing business owners to security trade-offs
\item Choosing a secure and well maintained framework
	%(more on frameworks later)
\item Devising a secure configuration for deployment
\end{itemize}
}
\end{frame}
%------------------------------------------------

\subsection{Penetration testing}

\begin{frame}
\frametitle{Penetration testing}
\includeimage[width=0.7\linewidth,height=0.7\textheight]{pentesting}
\note{
\textbf{Penetration testing} (or pentesting) is the action of attacking a
system with the consent of the owner in the goal of finding security holes.
\begin{itemize}
\item Identifying vulnerabilities
\item Exploit vulnerabilities
\item Testing ability to detect and respond to an attack
\item Social engineering
\end{itemize}
}
\end{frame}

\begin{frame}
\frametitle{Tools - Vulnerabilities detector}
\begin{columns}
\begin{column}{0.55\linewidth}
\begin{figure}
\includeimage[width=0.8\linewidth]{logo-metasploit}
{\scriptsize \url{http://www.metasploit.com/}}
\end{figure}
\end{column}
\begin{column}{0.45\linewidth}
\begin{figure}
\includeimage[width=0.8\linewidth]{logo-nessus}
{\scriptsize \url{http://www.tenable.com/}}
\end{figure}
\end{column}
\end{columns}
\note{
\begin{itemize}
\item \href{http://www.metasploit.com/about/penetration-testing-basics/}{Metasploit}:
	identifies vulnerabilities, automates assessments; attack a host
\item \href{http://www.tenable.com/products/nessus}{Nessus}:
	vulnerability scanner
\end{itemize}
}
\end{frame}

\begin{frame}
\frametitle{Tools - Network}
\begin{columns}
\begin{column}{0.5\linewidth}
\begin{figure}
\includeimage[width=0.8\linewidth]{logo-nmap}
NMAP
\newline {\scriptsize \url{http://nmap.org/}}
\end{figure}
\end{column}
\begin{column}{0.5\linewidth}
\begin{figure}
\includeimage[width=0.8\linewidth]{logo-webscarab}
WebScarab
\newline {\scriptsize \url{https://www.owasp.org/}}
\end{figure}
\end{column}
\end{columns}
\note{
\begin{itemize}
\item \href{http://nmap.org/}{nmap}:
	port scanning
\item \href{https://www.owasp.org/index.php/Webscarab}{WebScarab}:
	intercepting proxy
\end{itemize}
}
\end{frame}

\begin{frame}
\frametitle{Tools - All in one}

\includeimage[width=\linewidth,height=0.6\textheight]{bt5-screenshot}
\begin{center}
BackTrack
\newline {\scriptsize \url{http://www.backtrack-linux.org/}}
\end{center}
\note{
\href{http://www.backtrack-linux.org/}{BackTrack}:
	Linux distribution
}
\end{frame}

%------------------------------------------------
\begin{frame}
\frametitle{Maintenance}
\includeimage[width=0.7\linewidth,height=0.7\textheight]{patched}
\note{
\begin{block}{Patching}
When vulnerabilities have been discovered they have to be patched as soon as
possible.
Beware of regressions and new security breaches.
\end{block}
}
\end{frame}

%------------------------------------------------
\subsection{Using frameworks}

%------------------------------------------------
\begin{frame}
\frametitle{Frameworks}
\begin{quote}
\begin{center}
\LARGE
Your application is not the first.
\end{center}
\end{quote}
\note{
Frameworks ease development while providing security mechanisms:
\begin{itemize}
\item Authentication support
\item Escaping of values
\item Standard input validation
\end{itemize}
\begin{block}{Disadvantages}
\begin{itemize}
\item Changing a framework may prove difficult
\item Bug in a framework is known
	(version management needed)
\end{itemize}
\end{block}
}
\end{frame}

%------------------------------------------------
\begin{frame}
\frametitle{Frameworks - examples}

% TODO mettre une image avec plein de frameworks ESAPI, supprimer Spring

\begin{columns}
\begin{column}{0.55\linewidth}
\begin{figure}
\includeimage[width=0.8\linewidth]{logo-spring}
\scriptsize{\url{http://www.springsource.org/}}
\end{figure}
\begin{figure}
\includeimage[width=0.8\linewidth]{apache-shiro}
\scriptsize{\url{https://shiro.apache.org/}}
\end{figure}
\end{column}
\begin{column}{0.45\linewidth}
\begin{figure}
\includeimage[width=0.8\linewidth]{logo-esapi}
ESAPI
\newline \scriptsize{\url{https://www.owasp.org/}}
\end{figure}
\begin{figure}
\includeimage[width=0.8\linewidth]{dotnet-logo}
\scriptsize{\url{https://www.microsoft.com/net}}
\end{figure}
\end{column}
\end{columns}
\note{
\begin{exampleblock}{Spring}
Spring is one of the most used framework for developing web application using
Java. It does feature many components including some dedicated to security.
\end{exampleblock}
\begin{exampleblock}{ESAPI}
The OWASP Enterprise Security API is a free, open source, web application
security control library that is designed to make it easier for programmers
to retrofit security into existing applications.
The ESAPI libraries also serve as a solid foundation for new development.
\end{exampleblock}
}
\end{frame}

%------------------------------------------------
\subsection{Security in application servers}

%------------------------------------------------
\begin{frame}
\frametitle{Security in application servers}
\begin{block}{Application server}
The program that runs the web application
\end{block}
\vspace{1em}
\begin{itemize}
\item Authentication management
\item Directory listing disabling
\item User session management
\item Error handling
\item Input encoding
\end{itemize}
\note{
Application servers offer various services that can help you configure
web applications.
\begin{itemize}
\item Authentication management
\item Directory listing disabling
\item User session management
\item Error handling
\item Input encoding
\end{itemize}
}
\end{frame}

%------------------------------------------------
\subsection{Application Firewall}

%------------------------------------------------
\begin{frame}
\frametitle{Web Application Firewall}
\includeimage[width=0.7\linewidth,height=0.7\textheight]{firewall}
\note{
\textbf{Web Application Firewall} (WAF) are filters which inspects the traffic
that goes in and out your application and are able to block inappropriate
content.
\begin{block}{Virtual Patching}
Virtual Patching is a security policy layer that prevents the exploitation of a known
vulnerability until a patch is released. It does so by filtering the
traffic that goes to and from the application server and blocking
requests that are exploits.
\end{block}
}
\end{frame}

%------------------------------------------------
\subsection{HTTPS}

%------------------------------------------------
\begin{frame}
\frametitle{Secure protocols}
\includeimage[width=0.5\linewidth]{https}
\begin{columns}
\begin{column}{0.5\linewidth}
\centering\Huge\bfseries
SSL
\end{column}
\begin{column}{0.5\linewidth}
\centering\Huge\bfseries
TLS
\end{column}
\end{columns}
\note{
\textbf{TLS} (Transport Layer Security) is the successor of \textbf{SSL}
(Secured Socket Layer).
Equal in terms of security.
\begin{block}{What provides HTTPS?}
HTTPS provides \emph{encryption} and \emph{authentication} for HTTP.
\newline Protects against eavesdropping and spoofing.
\end{block}
\vspace{1em}
HTTPS is a part of a solution, it does not ensure full security.
\begin{block}{SSL vs TLS}
TLS is an open-community standard and allows both secure and unsecure
communications.
\newline TSL~1.0 is equivalent to SSL~3.1
\end{block}
}
\end{frame}

%------------------------------------------------
\subsection{Browser restrictions}

%------------------------------------------------
\begin{frame}
% TODO Image
% * Ne pas descendre Flash dans les explications ("il faut etre a jour")
% * "Javascript" => "JavaScript"
% * Ajouter le sandbox
% voir les notes du slide
\frametitle{Browser restrictions}
\includeimage[width=0.9\linewidth,height=0.9\textheight]{browser-security2}
\note{
\begin{itemize}
\item Cross-domain XMLHttpRequest (XHR) is disabled by default for other websites
	\\ Exceptions defined using \texttt{Access-Control}
\item Possibility of disabling JavaScript and plug-ins (Flash)
\item Private navigation mode
	\\ Does not store data on the disk.
\item Cookies access is limited by
	\begin{itemize}
	\item URL
	\item Timestamp
	\item Flags (Secure, HTTP-Only)
	\end{itemize}
\item SSL Certificates are validated
\item Blacklists of known rogue sites
\end{itemize}
}
\end{frame}

%------------------------------------------------
\subsection{E-mail}
%------------------------------------------------
\begin{frame}
\frametitle{E-mail}
\includeimage[width=0.7\linewidth,height=0.7\textheight]{spam}
\note{
E-mail is one of the most used application.
And it's plagued by spam and spoofing. But solutions exists:

\begin{block}{Sender Policy Framework (SPF)}
SPF validates e-mails by verifying sender's IP address against DNS records.
Administrators can define allowed hosts.
\end{block}

\begin{block}{Domain Keys Identified Mail (DKIM)}
DKIM associates domain names and email messages using a digital signature
added to the header of messages, a public key is added to the DNS record of the
domain for the recipient to read the message's header.
\end{block}
}
\end{frame}

%------------------------------------------------
\subsection{Variety of environments}

%------------------------------------------------
\begin{frame}
\frametitle{Variety of environments}
\includeimage[width=0.8\linewidth,height=0.8\textheight]{various-environments}
\note{
\begin{itemize}
\item Servers {\small (and their configuration)}
\item Networks
\item Clients
	\begin{itemize}
	\item Browsers
		\\ hacks {\small (example: IE6)} are just more code to test
	\item Mobile
	\item M2M
	\end{itemize}
\item Networks: proxy, caches, etc.
\item M2M: Machine to Machine
\end{itemize}
}
\end{frame}



\section{Detection and Analysis}

\subsection{Penetration testing}

\begin{frame}
\frametitle{Pentesting}
\includeimage[width=0.7\linewidth,height=0.7\textheight]{pentesting}
\note{
\textbf{Penetration testing} (or pentesting) is the action of attacking a
system with the consent of the owner in the goal of finding security holes.
\begin{itemize}
\item Identifying vulnerabilities
\item Exploit vulnerabilities
\item Testing ability to detect and respond to an attack
\item Social engineering
\end{itemize}
}
\end{frame}

\begin{frame}
\begin{center}
Avoiding this situation.
\end{center}
\includeimage[width=\linewidth,height=0.7\textheight]{fence-fail}
\end{frame}

\begin{frame}
\frametitle{Tools}
\begin{columns}
\begin{column}{0.35\linewidth}
\includeimage[width=\linewidth]{pentesting-tools}
\end{column}
\begin{column}{0.65\linewidth}
\includeimage[width=\linewidth,height=0.7\textheight]{bt5-screenshot}
\end{column}
\end{columns}
\note{
\begin{itemize}
\item \href{http://nmap.org/}{nmap}:
	port scanning
\item \href{http://www.metasploit.com/about/penetration-testing-basics/}{Metasploit}:
	identifies vulnerabilities, automates assessments
\item \href{http://www.tenable.com/products/nessus}{Nessus}:
	vulnerability scanner
\item \href{https://www.owasp.org/index.php/Webscarab}{WebScarab}:
	intercepting proxy
\item \href{http://www.backtrack-linux.org/}{BackTrack}:
	Linux distribution
\end{itemize}
}
\end{frame}

\subsection{Honeypot}

\begin{frame}
\frametitle{Honeypot}
\includeimage[width=0.8\linewidth,height=0.7\textheight]{honeypot}
\note{
A \textbf{honeypot} is a mocked website or web application that is exposed for
attackers who try to penetrate some system. A honeypot is a non-patched server
that is monitored for attacks and serves no other purpose.
\begin{exampleblock}{\href{https://www.projecthoneypot.org/about_us.php}{Project Honey Pot}}
Anti-spammer project.
Web sites publish some fake e-mails and spammer's crawlers sending e-mail to
such addresses are blocked.
\end{exampleblock}
}
\end{frame}

\subsection{NIDS}

\begin{frame}
\frametitle{NIDS}
\includeimage[width=0.8\linewidth,height=0.7\textheight]{logo-snort}
\note{
\textbf{Network Intrusion Detection System} (NIDS) is a system that scans
network activity and tries to detect attacks.
\begin{exampleblock}{\href{http://www.snort.org/}{Snort}}
A NIDS, and an IPS (Intrusion Prevention System), snort can detect attacks and
take actions accordingly.
\end{exampleblock}
}
\end{frame}

\begin{frame}
\frametitle{NIDS - detection types}
\begin{itemize}
\item Signature-based detection
\item Statistical anomaly-based detection
\item Stateful protocol analysis
\end{itemize}
\note{
\begin{block}{Signature-based detection}
Finds a pattern in packets on the network.
Detects SQL injections in HTTP parameters, XSS attacks, etc.
\end{block}
\begin{block}{Statistical anomaly-based detection}
Looks for peaks in traffic, connection coming from a single host and other
anomalies.
Detects stealth port scans, brute force attacks, etc.
\end{block}
\begin{block}{Stateful protocol analysis}
Follows a protocol conversation.
Detects SYN port scans, etc.
\end{block}
}
\end{frame}

\subsection{Post-mortem analysis}

\begin{frame}
\frametitle{After an attack}
\includeimage[width=0.8\linewidth,height=0.7\textheight]{investigation}
\note{
\begin{itemize}
\item \emph{When} did the attack happen
	\begin{itemize}
	\item Were we attacked?
	\item Is the attack still in progress?
	\end{itemize}
\item \emph{What} was compromised
	\begin{itemize}
	\item Data was compromised?
	\item Data was modified?
	\item What systems failed to work?
	\end{itemize}
\item \emph{How} the attacker succeeded
	\begin{itemize}
	\item How to fix the problem?
	\item Prevent future attacks?
	\end{itemize}
\end{itemize}
}
\end{frame}

\begin{frame}
\frametitle{How to know?}
\begin{itemize}
\item Traffic logs and graphs
\item Application logs
\end{itemize}
\pause
\begin{columns}
\begin{column}{0.5\textwidth}
\includeimage[width=\linewidth]{dos-load}
\end{column}
\begin{column}{0.5\textwidth}
\includeimage[width=\linewidth]{dos-bw}
\end{column}
\end{columns}
\note{
Read the logs and alerts.
\begin{exampleblock}{DoS}
Peak in server load and in input traffic.
Once the attack is effective, the outbound traffic drops.
\end{exampleblock}
}
\end{frame}

\begin{frame}
\frametitle{Analysis - tools}
\includeimage[width=0.8\linewidth,height=0.7\textheight]{analysis-tools}
\note{
\begin{itemize}
\item \href{http://www.ntop.org/products/ntop/}{ntop}: traffic monitoring
platform
\item \href{http://www.intersectalliance.com}{Snare}: a collection of tools to
audit log data and facilitate log analysis
\item \href{http://www.ossec.net/}{OSSEC}: open-source alternative; performs
log analysis, file integrity checking, policy monitoring, rootkit detection,
real-time alerting and active response
\item Web log analysis software
	\begin{itemize}
	\item \href{http://www.webalizer.org/}{Webalizer}
	\item \href{http://awstats.sourceforge.net/}{AWStats}
	\item \href{http://www.weblogexpert.com/}{WebLog Expert}
	\end{itemize}
\end{itemize}
}
\end{frame}

\subsubsection{Noise as a diversion}

\begin{frame}
\frametitle{Noise as a diversion}
\includeimage[width=0.8\linewidth,height=0.7\textheight]{train_cache}
\note{
When an attack happens, it can be masked by another attack which would serve
as a diversion from the real attack.
\begin{exampleblock}{DoS as a diversion}
DoS generates lot of traffic and can mask attacks as generated logs are huge.
\end{exampleblock}
}
\end{frame}

\subsection{Performances}

\begin{frame}
\frametitle{Performances}
\includeimage[width=0.8\linewidth,height=0.7\textheight]{performance}
\note{
\begin{itemize}
\item Slow response time and high load might be a symptom of an attack
\item Slow application implies higher loads and may not scale and crash during
	an attack
\end{itemize}
}
\end{frame}



%-----------------------------------------------------------------------------

\section{Conclusion}

\begin{frame}
\includeimage[height=0.7\textheight]{security-cloud}
\note{
	There it is. After a few reminders on cryptography and identity management, we showed you how easy it can be to exploit applications. And how small vulnerabilities can sum up and have impacts on your applications, on your users and finally degrade your image.

	After talking about some hints to fix, avoid or mitigate those threats, we took a look at which tools can help you detect and analyze attacks.

	We hope that you enjoyed this summary. We hope that you will look at your applications with its security in mind. Not in an extreme way, throwing money out the window for the sake of building a new digital Fort Knox. But rather with balance in mind and knowing the trade off you make.

	We will now have a small question and answer session. But if you don't have time to ask you questions or if you're a little shy of doing it in public, feel free to contact me on Twitter or by email. You can also download the slides of this presentation on ... where we will also publish a paper with the content in this presentation.
}
\end{frame}

\begin{frame}
\begin{flushright}
\Huge
Thanks
\end{flushright}
\vspace{1cm}
\includeimage[width=6cm]{hermes}
\note{
	Thank you.
}
\end{frame}

%\begin{frame}
%\frametitle{References}
%\begin{itemize}
%\item TODO
%\end{itemize}
%\end{frame}

%-----------------------------------------------------------------------------

\end{document}
